\section{Actual flat hypersurface fibres remain Cartier}

\begin{lemma}\label{lem:kn-nadel-principal-quotient-exact}
\lean{EnvelopeProperty.Vanishing.NadelBertiniFlatFibre.mulMap,EnvelopeProperty.Vanishing.NadelBertiniFlatFibre.quotientMap,EnvelopeProperty.Vanishing.NadelBertiniFlatFibre.mulMap_exact}
\leanok
For an algebra $B$ over a commutative ring $A$ and $h\in B$, the actual
$A$-linear maps $B\xrightarrow{h}B\to B/(h)$ form an exact pair, and
the actual quotient map is surjective.
\end{lemma}
\begin{proof}
The kernel of the quotient map is the original principal ideal $(h)$.
Membership in this ideal is precisely membership in the image of actual
multiplication by $h$.  The native quotient map is surjective.
\end{proof}

\begin{theorem}\label{thm:kn-nadel-flat-principal-fibre-regular}
\lean{EnvelopeProperty.Vanishing.NadelBertiniFlatFibre.mulMap_lTensor_injective,EnvelopeProperty.Vanishing.NadelBertiniFlatFibre.mulMap_lTensor_apply,EnvelopeProperty.Vanishing.NadelBertiniFlatFibre.tensor_regular}
\leanok
\uses{lem:kn-nadel-principal-quotient-exact}
Suppose $h$ is a nonzerodivisor in the actual $B$, and $B/(h)$ is flat
over $A$.  For every commutative $A$-algebra $S$, the actual element
$1\otimes h$ is a nonzerodivisor in $S\otimes_A B$.
No flatness assumption on $S$ is required.
\end{theorem}
\begin{proof}
The preceding pair is a short exact sequence.  The native flat-cokernel
tensor-exactness theorem implies that tensoring multiplication by $h$
with any $A$-module remains injective.  On each actual pure tensor
this map is $s\otimes b\mapsto s\otimes hb$, which equals multiplication
by $1\otimes h$ in the actual tensor algebra.  Additivity extends the
identity to every tensor.  Hence multiplication by that same element
is injective, proving the assertion.
\end{proof}

For the universal incidence equation, the normalized coefficient of one
parameter is one, so its actual equation is a nonzerodivisor before
specialization.  Restricting to the proved good parameter open makes the
actual incidence morphism flat.  The theorem therefore supplies the
genuine Cartier condition on each original field fibre.  This is the
missing extra argument beyond smoothness of a fibre alone; its eventual
application still retains the actual incidence charts and their quotient
identifications.
